\documentclass[12pt]{article}
\usepackage[T1]{fontenc}
\usepackage[utf8]{inputenc}
\usepackage{lmodern}
\usepackage{textcomp}
\usepackage{enumitem}
\usepackage{geometry}
\geometry{margin=1in}
\begin{document}
\begin{center}
Answer Key - Political Science - Undergraduate Level Assessment 10 \
\small (Answers are ordered exactly as in the question paper.)
\end{center}

\section*{Multiple Choice}
\begin{enumerate}[label=\arabic*.]
\item \textbf{Answer:} C
\\ \textit{Options:} A) Environmental activists; B) Labor unions; C) For-profit business; D) Religious institutions
\\ \textit{Note:} The for-profit business sector has created the largest number of political action committees since the 1970 s due to its significant financial resources and the desire to influence legislation and policy that affects their economic interests.
\item \textbf{Answer:} A
\\ \textit{Options:} A) Cautious support for Gorbachev and Yeltsin; B) Unreserved support for Gorbachev and Yeltsin; C) Denouncement of Gorbachev, but support for Yeltsin; D) Denouncement of Gorbachev and Yeltsin
\\ \textit{Note:} Under George H.W. Bush, U.S. relations with Russia were characterized by cautious support for leaders Mikhail Gorbachev and Boris Yeltsin as the U.S. aimed to encourage democratic reforms and stabilize the post-Cold War transition, while remaining mindful of the complexities of Russian politics.
\item \textbf{Answer:} A
\\ \textit{Options:} A) It criticized international organizations, rather than trying to strengthen them; B) It expanded NATO to include former Soviet states; C) It focused on a more personal style of leadership; D) It increased international support for the United States
\\ \textit{Note:} The correct answer reflects the George W. Bush administration's shift towards unilateralism in foreign policy, characterized by a skepticism of international organizations such as the United Nations, as seen in their actions regarding the Iraq War and their withdrawal from various international treaties.
\item \textbf{Answer:} C
\\ \textit{Options:} A) Having the capability to deter the most powerful rival; B) Having the capability to deter smaller states; C) How nuclear attacks are identified and responded to; who controls the weapons; D) Bureaucratic politics provides no information about nuclear proliferation and use
\\ \textit{Note:} The correct answer highlights the critical concerns of bureaucratic politics in nuclear weapons management, emphasizing the importance of decision-making processes regarding the identification and response to nuclear threats, as well as the control and authority over these weapons, which are essential for ensuring national and global security.
\item \textbf{Answer:} D
\\ \textit{Options:} A) Ideological differences; B) A power vacuum; C) The expansionist nature of the Soviet Union; D) Both b and c
\\ \textit{Note:} Realists argue that the onset of the Cold War can be attributed to the anarchic nature of the international system, which led to power struggles between the United States and the Soviet Union, as well as differing ideologies and national interests that heightened tensions and competition.
\end{enumerate}

\section*{True / False}
\begin{enumerate}[label=\arabic*.]
\setcounter{enumi}{5}
\item \textbf{Answer:} False
\\ \textit{Options:} A) True; B) False
\\ \textit{Note:} Krauthammer argued that the post-Cold War system was vulnerable to significant threats from rising powers and regional conflicts, indicating that dangers persisted despite the absence of a bipolar superpower rivalry.
\item \textbf{Answer:} False
\\ \textit{Options:} A) True; B) False
\\ \textit{Note:} The statement is false because the first Human Development Report was published in 1990, not 1987, which established the foundation for development studies.
\item \textbf{Answer:} True
\\ \textit{Options:} A) True; B) False
\\ \textit{Note:} In 'Federalist No. 10,' James Madison argues that a federal system, through its division of powers and representation, creates a multiplicity of interests and factions, making it challenging for any single faction to dominate and gain governing power.
\item \textbf{Answer:} False
\\ \textit{Options:} A) True; B) False
\\ \textit{Note:} Intergovernmental organizations typically lack the authority and mechanisms for robust enforcement, relying instead on member states' voluntary compliance and cooperation to adhere to international agreements.
\item \textbf{Answer:} True
\\ \textit{Options:} A) True; B) False
\\ \textit{Note:} The National Security Council (NSC) is indeed responsible for coordinating American foreign and military policy as it brings together key officials from various government agencies to advise the President on national security and foreign relations matters.
\end{enumerate}

\section*{Long Form Response}
\begin{enumerate}[label=\arabic*.]
\setcounter{enumi}{10}
\item \textbf{Answer:} Security studies in the contemporary world face several significant challenges, including the rise of non-state actors, the complexity of cybersecurity threats, and the impact of climate change on global security dynamics. The proliferation of technology has altered traditional security paradigms, creating new vulnerabilities and necessitating a reevaluation of strategies. Additionally, the interconnectedness of global politics means that issues such as economic instability and health crises can directly affect national and international security. These challenges will shape the future of the field by pushing scholars and practitioners to adopt interdisciplinary approaches and develop more adaptive frameworks that account for an ever-evolving security landscape. Ultimately, addressing these challenges will require innovative thinking and collaboration across various domains within and beyond security studies.
\\ \textit{Note:} The answer correctly identifies key contemporary challenges in security studies, such as non-state actors, cybersecurity threats, and climate change, while emphasizing the need for a reevaluation of traditional security strategies in light of global interconnectedness and emerging vulnerabilities.
\item \textbf{Answer:} The use of military force to expel social groups from a state has profound implications for national identity, human rights, and international relations. Such actions often lead to a fractured national identity, as the state may redefine itself through exclusion rather than inclusion, fostering divisions among its populace. Human rights violations become a significant concern, as military interventions typically result in displacement, violence, and loss of life, undermining the fundamental rights of affected individuals. Furthermore, these actions can strain international relations, as states that engage in such practices may face condemnation, sanctions, or intervention from the global community, leading to isolation or conflict. Ultimately, the use of military force in this manner raises ethical questions and challenges the legitimacy of the state's governance.
\\ \textit{Note:} The answer correctly highlights that the use of military force to expel social groups can fracture national identity by promoting exclusion, while also leading to severe human rights violations that undermine fundamental rights, and strains international relations due to global condemnation and potential violations of international law.
\end{enumerate}

\vfill
\noindent\textit{End of Answer Key}
\end{document}